\paragraph{A completely specified complete-intersection realization.}
The following finite formula fixes the degree, every section, and every
coefficient. Use the 24 coordinates $x_{r,n,a,b}$, where $1\leq r\leq3$
and $n,a,b\in\{0,1\}$. The central $\mu_2$ has weight $n$, and
\[
 A(x_{r,n,a,b})=x_{r,n,a+1,b},\qquad
 B(x_{r,n,a,b})=(-1)^{na}x_{r,n,a,b+1}.                    \tag{E4.H.4}
\]
The indices $a,b$ are modulo two. These are three copies of the regular
coordinate representation, so $P=\mathbf P^{23}$ is the projectivization
of the threefold direct sum of the contragredient regular representation
in Raynaud's Lemma~4.2.2. Its free-action open $U$ has complement of
fibrewise codimension at least three.

Order the coordinates lexicographically. Let $\mathcal A$ be the ordered
finite set of degree-192 monomials with even total $\mu_2$ weight. Form
\[
 F=\sum_{\alpha\in\mathcal A}u_\alpha x^\alpha,
 \qquad F\,A(F)\,B(F)\,AB(F)=\sum_{|\beta|=4}u^\beta p_\beta(x).
                                                               \tag{E4.H.5}
\]
The actions fix the $u$ variables. On even-weight forms $A,B$ commute,
so every integral polynomial $p_\beta$ is $G$-invariant of degree 768.
Retain zero polynomials and repetitions; no basis choice is involved.
Let $\mathcal B$ be the ordered set of the multiindices $\beta$ above.
Define
\[
 k=\overline{\mathbf F_2(c_{i\beta}:1\leq i\leq21,
                                  \ \beta\in\mathcal B)},
 \quad R=W(k),\qquad
 s_i=\sum_{\beta\in\mathcal B}[c_{i\beta}]p_\beta.          \tag{E4.H.6}
\]
Here all variables are algebraically independent and brackets are
Teichm\"uller representatives. These are finite definitions, with
\[
 h=|\mathcal A|=\tfrac12\{\binom{215}{23}+\binom{107}{11}\},
 \qquad |\mathcal B|=\binom{h+3}{4}.
\]

To type the descent precisely, put $M=P/\mu_2$ and $Q=M/H=P/G$, where
$H=(\mathbf Z/2)^2$. The central norm of $\mathcal O_P(1)$ is a bundle
$\mathcal N$ on $M$ with pullback $\mathcal O_P(2)$, naturally
$H$-linearized. Put $\mathcal L=\operatorname{Norm}_H(\mathcal N)$ on
$Q$. It is ample with pullback $\mathcal O_P(8)$. These are Raynaud's
norm bundles with their factorization made explicit
\cite[proof of Proposition~4.2.3, p.~137]{BCERaynaud1979}.
Even-weight degree-192 forms are sections of $\mathcal N^{96}$.
Their $H$-norms, hence all $p_\beta$, are sections of $\mathcal L^{96}$.
These integral formulas avoid any assumption about reduction of a chosen
basis of sections. Define
\[
 X=V(s_1,\ldots,s_{21})\subset Q,
 \qquad Y=q^{-1}(X)\subset P.                              \tag{E4.H.7}
\]
Thus $Y$ is specified by 21 equations of the fixed degree 768.

\paragraph{Why this fixed degree suffices.}
Write $V=U/G$. The maps $U\to U/\mu_2\to V$ are respectively a degree-two
$\mu_2$-torsor and a degree-four finite \emph{\'etale} $H$-torsor.
The scheme $V$ is smooth of relative dimension 23: flatness descends
from $U$, and geometric regularity descends through the faithfully flat
finite maps on fibres. The complement $Q_k\setminus V_k$ has dimension
at most 20.

The sections $p_\beta$ generate $\mathcal L^{96}$ on $Q_k$. Indeed even
coordinate powers generate $\mathcal N^{96}$ on $M_k$. Since $k$ is
infinite, a linear combination can be chosen nonzero at all points of
any finite $H$-orbit. Its norm is nonzero at the corresponding point of
$Q_k$, and every norm is a linear combination of the coefficient
sections $p_\beta$.

They also span every first jet at $v\in V(k)$. The first infinitesimal
neighborhood of $v$ has length 24; its inverse image $Z$ in $P_k$ has
length $8\cdot24=192$. A length-$\ell$ subscheme of projective space
over an infinite field is interpolated by $\mathcal O(d)$ for
$d\geq\ell-1$: choose a hyperplane avoiding its support and let $W_d$
be the span of degree-at-most-$d$ monomials in its length-$\ell$
coordinate algebra. If $W_{d+1}=W_d$, this space is stable under every
coordinate and is the whole algebra. Otherwise its dimension strictly
increases. Since $\dim W_0=1$, it reaches the algebra by $d=\ell-1$.
Thus degree-192 forms interpolate all sections on $Z$.

Taking $\mu_2$ invariants preserves that surjection: projection onto
weight zero is exact even in characteristic two. Sections of
$\mathcal N^{96}$ therefore interpolate arbitrary first jets at the
four points above $v$ in $U/\mu_2$. Prescribe the desired jet at one
point and unit jets at the other three, in $H$-compatible
trivializations. The norm has the desired jet at $v$. Consequently
the span of the $p_\beta$ surjects onto every first jet on $V_k$.

The incidence proof of Bertini now makes 21 generic sections smooth on
$V_k$, while global generation and the dimension bound 20 make them
avoid its complement. The independent coefficients in (E4.H.6) specify
such a generic tuple. Amplitude makes the intersection a nonempty
surface. Thus $X_k$ is smooth and contained in $V_k$. Since $Q$ is smooth
near $X_k$ and the 21 equations cut its smooth special fibre in codimension
21, they form a fibrewise regular sequence there. The local flatness
criterion makes $X/R$ flat near $X_k$, and the smooth special fibre then
makes $X/R$ smooth on that neighborhood. These properties hold throughout
$X$: their bad loci are closed in the projective
$R$-scheme $X$, and a nonempty closed image in the local base must meet
the closed point. Therefore $X/R$ is smooth projective, $Y\to X$ is a
$G$-torsor, and the 21 pullback equations are a relative complete
intersection in $P^{23}_R$. This verifies all hypotheses of Raynaud's
Proposition~4.2.3 with an explicit power $m=96$.

A complete intersection of dimension at least two has
$H^0(\mathcal O)=k$, $H^1(\mathcal O)=0$, and trivial Picard tau-group,
including when singular or nonreduced; the last statement is
\cite[Theorem~1.1.8(iii), proved in \S3.4]{BCECotnerZavyalov2024}.
Hence $Y/R$ has universal $H^0=\mathcal O_R$ and
$\operatorname{Pic}^\tau_{Y/R}=0$, by the fibrewise assertion and the
unramified identity-section criterion. Geometric connectedness of $Y$
and its surjection onto $X$ make the smooth fibres of $X$ geometrically
integral and normal. The coordinate-check script verifies the signed
representation and the finite index cardinalities. It does not expand
the enormous norm polynomial or claim a computational smoothness test;
the interpolation and incidence arguments supply that proof.

